\chapter{Laser Source Classes: Noise and Practical Notes}
\label{appendix:laserSourceNotes}
% Referenced from: Table 3.3 caption (\autoref{tab:laser-sources-notes}).
% Needs the array and booktabs packages (already loaded for Table 3.1).
\begin{table}[H]
\centering\small
\renewcommand{\arraystretch}{1.2}
\begin{tabular}{@{}>{\raggedright\arraybackslash}p{0.22\textwidth}>{\raggedright\arraybackslash}p{0.76\textwidth}@{}}
\toprule
\textbf{Source type} & \textbf{Noise and practical notes} \\
\midrule
Fabry--P\'erot laser diode &
Oscillates on several longitudinal modes simultaneously, so the output is not single-frequency and cannot be mixed coherently with a local oscillator; unsuitable for coherent detection \cite{harris_role_1998}. \\
\addlinespace
DFB laser diode &
Single-frequency, but with a linewidth of the order of MHz, which gives rise to significant PIIN in a CW system unless the LO and transmit path lengths are closely matched (\autoref{sec:testingPIIN}). Relaxation-oscillation peak at $\sim$\SI{1}{\giga\hertz}, far above the Doppler band, so RIN is flat and low in the frequency band of interest \cite{cariou2006laser}. Supplied as a bare butterfly module; needs an external current driver and temperature controller. Selected for this project. \\
\addlinespace
Planar external-cavity semiconductor &
kHz-range linewidth, so PIIN and Doppler-peak broadening are greatly reduced compared to a DFB diode. Supplied with an integrated driver, and marketed specifically for wind lidar. \\
\addlinespace
Er-doped DFB fibre laser &
The common source of choice for CW anemometry owing to its very narrow linewidth \cite{cariou2006laser,hill_remote_nodate,abdelazim_development_2015}. Relaxation-oscillation peak near \SI{1}{\mega\hertz}, however, so RIN is higher in the frequency band of interest than for a DFB diode \cite{cariou2006laser}. \\
\bottomrule
\end{tabular}
\caption{Noise and practical notes for the laser source classes compared in \autoref{tab:design-laser-sources}.}
\label{tab:laser-sources-notes}
\end{table}
